% Generated draft: requires mathematical and source review.
\documentclass[11pt]{article}
\usepackage[margin=25mm]{geometry}
\usepackage{fontspec}
\setmainfont{Latin Modern Roman}
\usepackage{amsmath,amssymb,mathtools,unicode-math}
\setmathfont{Latin Modern Math}
\usepackage{enumitem,needspace,xurl,hyperref}
\setlength{\parindent}{0pt}
\setlength{\parskip}{5pt}
\setlength{\emergencystretch}{3em}
\newcommand{\N}{\mathbb{N}}
\newcommand{\Z}{\mathbb{Z}}
\newcommand{\Q}{\mathbb{Q}}
\newcommand{\R}{\mathbb{R}}
\newcommand{\C}{\mathbb{C}}
\newcommand{\F}{\mathbb{F}}
\newcommand{\A}{\mathbb{A}}
\newcommand{\PP}{\mathbb{P}}
\newcommand{\E}{\mathbb{E}}
\newcommand{\eps}{\varepsilon}
\newcommand{\dd}{\,\mathrm{d}}
\begin{document}
\section*{UnsolvedMath ID: 800988}
\textbf{Source problem code:} catalogue entry
\subsection{Definitions and mathematical statement}
Let \(M\) be a smooth manifold of dimension \(n\geq 3\). In the usual smooth, strongly convex convention, a Finsler metric is a function \(F:TM\to[0,\infty)\), continuous on \(TM\) and smooth away from the zero section, such that \(F(x,v)>0\) for \(v\neq0\), \(F(x,tv)=tF(x,v)\) for \(t>0\), and the fibre Hessian of \(\tfrac12F^2\) is positive definite at every nonzero tangent vector. Reversibility, meaning \(F(x,-v)=F(x,v)\), is not imposed in this convention; the supplied source does not specify whether it intends any reversibility restriction.

Two Finsler metrics are projectively equivalent if they have the same unparameterized geodesics, subject to the unresolved orientation convention noted below. Given a smooth closed one-form \(\beta\) on an open set \(U\subseteq M\), a trivial projective change is
\[
F_\beta(x,v)=F(x,v)+\beta_x(v),\qquad (x,v)\in TU,
\]
provided this expression is again an admissible Finsler metric. The supplied source states that every sufficiently small closed one-form gives such a change without changing the geodesics; it attributes the terminology to \([\mathrm{Matveev2013}]\).

\textbf{Question (Burns--Matveev, Question 69).} In dimension \(n\geq3\), is it true that, for a generic Finsler metric \(F\), every Finsler metric \(\widetilde F\) projectively equivalent to \(F\) is locally proportional to a trivial projective change of \(F\)? Reading ``locally proportional'' as multiplication by a positive constant on each neighborhood, the proposed local conclusion is that, for every \(p\in M\), there exist an open neighborhood \(U\ni p\), a constant \(c_U>0\), and a smooth one-form \(\beta_U\) on \(U\) such that
\[
\dd\beta_U=0,\qquad
\widetilde F(x,v)=c_U\bigl(F(x,v)+\beta_{U,x}(v)\bigr)
\quad\text{for every }(x,v)\in TU,
\]
with \(F+\beta_U\) an admissible Finsler metric.

\textbf{Unresolved qualifications.} The source does not specify the space of metrics or topology underlying ``generic,'' the regularity or reversibility class, whether projective equivalence preserves the orientation of geodesics, or whether the metrics and equivalence are initially local or global. It also does not explicitly define ``locally proportional''; the displayed formula records the constant-factor interpretation. Thus a fully quantified genericity assertion requires clarification of these choices.
\subsection{Short English statement}
The supplied question asks whether, in dimension at least three, a generic Finsler metric is determined locally by its unparameterized geodesics up to positive scaling and addition of a closed one-form. The meanings of genericity and several metric and equivalence conventions remain unspecified. The registry labels the question open; this draft does not independently verify that status.
\subsection{Sources}
\begin{itemize}
\item Supplied local UnsolvedMath registry record, permanent ID \(800988\), source problem code \texttt{catalogue entry}, titled \emph{Open Problems and Questions about Geodesics --- Question 69}. The record attributes the problem to Burns and Matveev, gives the year 2019, and labels its status \texttt{open}.
\item Burns and Matveev, \emph{Open Problems and Questions about Geodesics} (2019, as identified by the supplied record), Question 69. Supplied source URL: \url{https://arxiv.org/abs/1308.5417}. The provenance records extraction by \texttt{arxiv\_latex} and an access date of 2026-10-16. These are supplied provenance statements, not an independent source visit or literature verification.
\item The extracted statement cites \([\mathrm{Matveev2013}]\) for the terminology ``trivial projective change.'' No full bibliographic entry or separate URL for this reference is supplied, so its precise identity is unresolved.
\end{itemize}
\end{document}
